\section{Original affine local-map readback and restriction}

Let $W$ be an original affine open in a scheme $Z$, and let $M,N$ be
original module sheaves. Assume the restriction of $M$ to the actual
$\operatorname{Spec}\Gamma(Z,W)$ is quasicoherent. All comparisons below
are the native restriction functors, their original identity/composition
isomorphisms, and the original affine associated-sheaf adjunction.

\begin{definition}[The original affine map of a section morphism]
\label{def:di-canonical-affine-local-map}
\leanok
\lean{EnvelopeProperty.DI.AffineLocalMapRestriction.specHomOfSections,EnvelopeProperty.DI.AffineLocalMapRestriction.specHomOfSections_zero,EnvelopeProperty.DI.AffineLocalMapRestriction.specHomOfSections_add,EnvelopeProperty.DI.AffineLocalMapRestriction.specHomOfSections_sub,EnvelopeProperty.DI.AffineLocalMapRestriction.affineSpecMap,EnvelopeProperty.DI.AffineLocalMapRestriction.chartRestrictionIso,EnvelopeProperty.DI.AffineLocalMapRestriction.affineLocalMap,EnvelopeProperty.DI.AffineLocalMapRestriction.affineLocalMap_zero,EnvelopeProperty.DI.AffineLocalMapRestriction.affineLocalMap_add,EnvelopeProperty.DI.AffineLocalMapRestriction.affineLocalMap_sub,EnvelopeProperty.DI.AffineLocalMapRestriction.affineLocalMap_postcompose}
A genuine $\Gamma(Z,W)$-linear map on original sections determines a map
$M|_W\to N|_W$: first conjugate by the original affine section units,
then use the native associated-sheaf adjunction and its invertible source
counit, and finally transport through the original affine-chart isomorphism.
\end{definition}
\begin{proof}
The source's genuine quasicoherence makes its original associated-sheaf
counit invertible. Compose its inverse with the native adjoint of the
original section map. The native adjunction triangle gives exact readback
on sections over the spectrum. Restriction along the original affine-chart
isomorphism and the original restriction composition comparisons produces
the stated map on the original open scheme. No morphism or comparison
certificate is supplied besides the original section morphism. The same
native additive adjunction preserves zero, addition and subtraction;
source-counit naturality gives compatibility with original global
postcomposition.
\end{proof}

\begin{lemma}[Whole original spectrum readback]
\label{lem:di-affine-local-map-spec-readback}
\leanok
\lean{EnvelopeProperty.DI.AffineLocalMapRestriction.specHomOfSections_map,EnvelopeProperty.DI.AffineLocalMapRestriction.affineSpecMap_sections,EnvelopeProperty.DI.AffineLocalMapRestriction.moduleSections_naturality,EnvelopeProperty.DI.AffineLocalMapRestriction.chartRestriction_naturality,EnvelopeProperty.DI.AffineLocalMapRestriction.specReturn,EnvelopeProperty.DI.AffineLocalMapRestriction.specReadbackIso,EnvelopeProperty.DI.AffineLocalMapRestriction.affineLocalMap_spec_readback,EnvelopeProperty.DI.AffineLocalMapRestriction.sectionReadback,EnvelopeProperty.DI.AffineLocalMapRestriction.sectionReadback_affineLocalMap,EnvelopeProperty.DI.AffineLocalMapRestriction.specReadback_naturality}
\uses{def:di-canonical-affine-local-map}
Restricting this same local map back along the inverse original chart
isomorphism, then conjugating by the original restriction identity and
composition isomorphisms, returns exactly the original spectrum map.
Consequently its native spectrum-section readback, conjugated by the
original affine section units, is exactly the original section morphism.
\end{lemma}
\begin{proof}
Write $F$ for restriction along the original chart isomorphism and $G$
for restriction along its inverse. The original composition and identity
comparisons give a natural isomorphism $FG\cong1$ on spectrum modules.
The local map is the conjugate of $F$ applied to the spectrum map by the
original chart restriction comparisons. Applying $G$ and conjugating by
the corresponding comparisons cancels those same original isomorphisms.
Naturality of $FG\cong1$ then gives exactly the original spectrum map.
Apply the native section functor and the original affine section units;
the established adjunction triangle readback gives precisely
the supplied original section morphism.
\end{proof}

\begin{lemma}[An affine detects adjoints supported on a smaller open]
\label{lem:di-affine-adjoint-section-faithfulness}
\leanok
\lean{EnvelopeProperty.DI.AffineLocalMapRestriction.specMap_ext,EnvelopeProperty.DI.AffineLocalMapRestriction.localMap_ext_of_adjoint_at_affine,EnvelopeProperty.DI.AffineLocalMapRestriction.localMap_ext_of_adjoint_sections}
\uses{lem:di-affine-local-map-spec-readback}
Let $V\subset W$ with $W$ affine. Suppose $M$ restricted to the original
spectrum of $W$ is quasicoherent. Two maps $M|_V\to N|_V$ are equal if
their original restriction-adjoints have equal section maps on $W$.
No quasicoherence of the pushed target is required.
\end{lemma}
\begin{proof}
The original affine section comparison turns equality on $W$ into equality
of section maps after restriction to its original spectrum. The source
counit is invertible by its genuine quasicoherence, so the native affine
associated-sheaf adjunction makes this section functor faithful on maps
from that source. Thus the two adjoints agree after spectrum restriction.
The original inclusion of $V$ factors through that same spectrum by the
canonical open inclusion and original chart isomorphism. The native
restriction composition comparison therefore gives equality after
restriction to $V$. Under the original restriction--pushforward adjunction,
a map into the original pushforward from $V$ is determined by that
restriction followed by its original counit. Injectivity of the original
adjunction equivalence gives equality of the adjoints, and then equality
of the two original local maps.
\end{proof}

\begin{lemma}[Restriction naturality from original section naturality]
\label{lem:di-affine-local-map-restriction}
\leanok
\lean{EnvelopeProperty.DI.AffineLocalMapRestriction.sectionReadback_restrict,EnvelopeProperty.DI.AffineLocalMapRestriction.sectionReadback_comp,EnvelopeProperty.DI.AffineLocalMapRestriction.restrict_localAdjoint,EnvelopeProperty.DI.AffineLocalMapRestriction.localAdjoint_sections,EnvelopeProperty.DI.AffineLocalMapRestriction.affineLocalMap_restrict}
\uses{lem:di-affine-local-map-spec-readback,lem:di-affine-adjoint-section-faithfulness}
For original affine opens $V\subset W$, original linear section morphisms
which commute with the actual original section restrictions induce the
same actual restricted local sheaf map on $V$. No matching sheaf-map
hypothesis is needed.
\end{lemma}
\begin{proof}
Use the original restriction adjunction on the larger affine open.
Naturality of the original inverse-chart comparisons makes readback of a
restricted global map its original section map; cancellation of those same
comparisons makes readback preserve composition. The original adjunction
triangle therefore identifies a local map's same-open adjoint section map
with its exact readback followed by the original target unit. Apply this
formula to both affine constructions. Naturality of the original target
unit and of the smaller-open adjoint reduces equality on the containing
affine to the supplied section square, after composition with restriction
of the original pushed target. That restriction is invertible: both cut
opens are the same smaller open. Cancel it. The assumed square therefore
makes the two adjoint section maps equal. The source is genuinely
quasicoherent on that affine, so the proved affine adjoint faithfulness
makes the whole original adjoints equal. The original adjunction then
gives equality of the local maps on the smaller open. The original
factor-map equation for native local restriction identifies the first of
those local maps with the exact native restriction of the larger one.
Thus only the original section square is used; no equality of the desired
sheaf maps is assumed.
\end{proof}
